%%%%%%%%%%%%%%%%%%%%%%%%%%%%%%%%%%%%
% Tables
%%%%%%%%%%%%%%%%%%%%%%%%%%%%%%%%%%%%
% =========================
% Sum stat tables
% =========================

\begin{landscape}
\input{tab01-SumStat}
\end{landscape}

\newpage
\pagebreak

\input{tab53-SumStatOverAll-adults}

\newpage
\pagebreak
% =========================
% Identity by proxy
% =========================

\input{maineffect-proxy}

\pagebreak
\newpage

% =========================
% Number of observations by generation
% =========================

\input{tab02-observations-by-gen}

\newpage
\pagebreak

\input{tab54-observations-by-gen-adults}

\pagebreak
\newpage


% =========================
% Migration table
% =========================

\input{tab46-iat-migration-tab.tex}

\pagebreak
\newpage

% =========================
% Marriage table
% =========================


\input{tab41a-marriage}

\pagebreak
\newpage

% =========================
% Between-Model Tests of AME
% =========================

\begin{table}[H]
\centering
\caption{Between-Model Tests of Average Marginal Effects: Second-Generation AW vs.\ WA \label{tab:between-model-second}}
\centering
\fontsize{10}{12}\selectfont
\begin{threeparttable}
\begin{tabular}[t]{lccc}
\toprule
 & \multicolumn{1}{c}{Asian Only} & \multicolumn{1}{c}{White Only} & \multicolumn{1}{c}{Asian \& White} \\
\cmidrule(l{3pt}r{3pt}){2-2} \cmidrule(l{3pt}r{3pt}){3-3} \cmidrule(l{3pt}r{3pt}){4-4}
\midrule
\multicolumn{4}{l}{\textit{Panel A: AW vs.\ WA}} \\
\addlinespace
Anti-Asian Bias & $-0.024$ & $0.106$** & $-0.082$\\
 & ($0.055$) & ($0.044$) & ($0.062$)\\
\addlinespace
Female & $-0.056$* & $-0.044$ & $0.097$***\\
 & ($0.031$) & ($0.029$) & ($0.028$)\\
\addlinespace
College Graduate: Mother & $0.084$** & $-0.106$** & $0.021$\\
 & ($0.036$) & ($0.050$) & ($0.048$)\\
\addlinespace
College Graduate: Father & $-0.061$** & $0.069$** & $-0.002$\\
 & ($0.027$) & ($0.031$) & ($0.037$)\\
\bottomrule
\multicolumn{4}{l}{\rule{0pt}{1em}* p $<$ 0.1, ** p $<$ 0.05, *** p $<$ 0.01}\\
\end{tabular}
\begin{tablenotes}
\small
\item[1] \footnotesize{Each cell reports the difference in average marginal effects (AME) between the Asian father--White mother (AW) and White father--Asian mother (WA) models, estimated from equation (\ref{eq:multinomial_logit}) using multinomial logit. The difference is computed as $\widehat{\text{AME}}_{\text{AW}} - \widehat{\text{AME}}_{\text{WA}}$.}
\item[2] \footnotesize{Standard errors of the difference are computed from the bootstrap distributions of the two models. $z$-statistics test the null hypothesis that the marginal effects are equal across models.}
\item[3] \footnotesize{I include region $\times$ year fixed effects with controls for quartic age and local Asian population share. Each estimate is the median across 1{,}000 bootstrap resamples, with 95\% percentile confidence intervals.}
\end{tablenotes}
\end{threeparttable}
\end{table}


\pagebreak
\newpage

\begin{table}[H]
\centering
\caption{Between-Model Tests of Average Marginal Effects: Third-Generation by Asian Grandparent Count \label{tab:between-model-third}}
\centering
\fontsize{10}{12}\selectfont
\begin{threeparttable}
\begin{tabular}[t]{lccc}
\toprule
 & \multicolumn{1}{c}{Asian Only} & \multicolumn{1}{c}{White Only} & \multicolumn{1}{c}{Asian \& White} \\
\cmidrule(l{3pt}r{3pt}){2-2} \cmidrule(l{3pt}r{3pt}){3-3} \cmidrule(l{3pt}r{3pt}){4-4}
\midrule
\multicolumn{4}{l}{\textit{Panel A: One Asian GP vs.\ Two Asian GP}} \\
\addlinespace
Anti-Asian Bias & $0.102$* & $-0.059$ & $-0.047$\\
 & ($0.059$) & ($0.077$) & ($0.084$)\\
\addlinespace
Female & $-0.010$ & $-0.073$ & $0.088$*\\
 & ($0.031$) & ($0.048$) & ($0.050$)\\
\addlinespace
College Graduate: Mother & $-0.048$ & $-0.106$** & $0.160$***\\
 & ($0.032$) & ($0.048$) & ($0.052$)\\
\addlinespace
College Graduate: Father & $-0.021$ & $0.060$ & $-0.041$\\
 & ($0.031$) & ($0.058$) & ($0.061$)\\
\midrule
\multicolumn{4}{l}{\textit{Panel B: One Asian GP vs.\ Three Asian GP}} \\
\addlinespace
Anti-Asian Bias & $0.201$ & $-0.203$*** & $-0.028$\\
 & ($0.152$) & ($0.071$) & ($0.149$)\\
\addlinespace
Female & $0.031$ & $-0.118$*** & $0.087$**\\
 & ($0.031$) & ($0.043$) & ($0.040$)\\
\addlinespace
College Graduate: Mother & $-0.318$*** & $0.198$*** & $0.122$**\\
 & ($0.046$) & ($0.047$) & ($0.051$)\\
\addlinespace
College Graduate: Father & $0.293$*** & $-0.345$*** & $0.057$\\
 & ($0.049$) & ($0.065$) & ($0.062$)\\
\midrule
\multicolumn{4}{l}{\textit{Panel C: Two Asian GP vs.\ Three Asian GP}} \\
\addlinespace
Anti-Asian Bias & $0.098$ & $-0.145$*** & $0.019$\\
 & ($0.156$) & ($0.054$) & ($0.142$)\\
\addlinespace
Female & $0.042$ & $-0.045$ & $-0.000$\\
 & ($0.041$) & ($0.047$) & ($0.045$)\\
\addlinespace
College Graduate: Mother & $-0.270$*** & $0.304$*** & $-0.039$\\
 & ($0.054$) & ($0.055$) & ($0.059$)\\
\addlinespace
College Graduate: Father & $0.314$*** & $-0.405$*** & $0.098$\\
 & ($0.056$) & ($0.060$) & ($0.067$)\\
\bottomrule
\multicolumn{4}{l}{\rule{0pt}{1em}* p $<$ 0.1, ** p $<$ 0.05, *** p $<$ 0.01}\\
\end{tabular}
\begin{tablenotes}
\small
\item[1] \footnotesize{Each cell reports the difference in average marginal effects (AME) between the indicated models, estimated from equation (\ref{eq:multinomial_logit}) using multinomial logit. The difference is computed as $\widehat{\text{AME}}_{\text{Model A}} - \widehat{\text{AME}}_{\text{Model B}}$.}
\item[2] \footnotesize{Standard errors of the difference are computed from the bootstrap distributions of the two models. $z$-statistics test the null hypothesis that the marginal effects are equal across models.}
\item[3] \footnotesize{I include region $\times$ year fixed effects with controls for quartic age and local Asian population share. Each estimate is the median across 1{,}000 bootstrap resamples, with 95\% percentile confidence intervals.}
\item[4] \footnotesize{GP = Grandparent(s). Third-generation individuals are native-born with at least one grandparent born in an Asian country.}
\end{tablenotes}
\end{threeparttable}
\end{table}


\pagebreak
\newpage

% =========================
% Diagram of how
% generations are defined
% =========================

\pagebreak
\newpage

\begin{center}
\begin{figure}[H]
\caption{Diagram of the Three Different Generations of Asian Immigrants.}
\includegraphics[width=\textwidth, height=9cm]{diag.png} 
\label{fig:diag}
\end{figure}
\hfill%
\end{center}

% =========================
% Bias and Self-reported 
% Asian Identity in the Least 
% and Most Biased Places
% =========================


\begin{center}
\begin{figure}[H]
\caption{Bias and Self-reported Asian Identity in the Least and Most Biased Places} \label{fig:toptwobias}

% first
\begin{subfigure}{.9\textwidth}
\caption{Bias Over Time}
\centering
\includegraphics[width=.9\linewidth]{Bias_twostates.png} 
\label{fig:skiniat}
\end{subfigure}
% Second
\begin{subfigure}{.9\textwidth}
\caption{Self-reported Asian Identity Over Time}
\centering
\includegraphics[width=.9\linewidth]{Bias_twostates-asian.png} 
\label{fig:Asian-twostates}
\end{subfigure}
\caption*{\footnotesize{These two panels show the trends in bias (panel a) and self-reported Asian identity among Asian immigrants (panel b) of the least and most biased places in the data. The District of Colombia is the least biased geographical area, and North Dakota is the most biased. The bias units are in standard deviations. Self-reported Asian identity is among first, second, and third-generation Asian immigrants aged 17 and younger still living in intact families.\\
Bias data is from the 2004--2021 Harvard's Project Implicit Association Test scores, American National Election Studies (ANES), and state-level hate crimes against Asians. Identity data is from the 2004--2021 Current Population Survey (CPS).}}
\end{figure}
\end{center}


\pagebreak
\newpage

% =========================
% Anti-Asian Bias Maps
% =========================


\begin{center}
\begin{figure}[H]
\caption{Maps of State-level Association Test Bias Over Time Measure with Census Division Regional Boundaries}
\label{fig:skiniat-maps}
% first
\begin{subfigure}{.45\textwidth}
\caption{State-level Bias in 2004}
\centering
\includegraphics[width=0.9\linewidth]{2004skinmap.png} 
\label{fig:skiniat-map-2004}
\end{subfigure}
\hfill%
% Second
\begin{subfigure}{.45\textwidth}
\caption{State-level Bias in 2008}
\centering
\includegraphics[width=0.9\linewidth]{2008skinmap.png} 
\label{fig:skiniat-map-2006}
\end{subfigure}
\hfill%
% third
\begin{subfigure}{.45\textwidth}
\caption{State-level Bias in 2012}
\centering
\includegraphics[width=0.9\linewidth]{2012skinmap.png} 
\label{fig:skiniat-map-2008}
\end{subfigure}
\hfill%
% fourth
\begin{subfigure}{.45\textwidth}
\caption{State-level Bias in 2016}
\centering
\includegraphics[width=0.9\linewidth]{2016skinmap.png} 
\label{fig:skiniat-map-2010}
\end{subfigure}

\caption*{\footnotesize{This figure shows the state-level bias index in different years in the sample. The bias units are in standard deviations and ranges from low to high bias. Bias index is constructed following \textcite{lubotskyInterpretationRegressionsMultiple2006}. The data is from the 2004--2021 Harvard's Project Implicit Association Test scores, American National Election Studies (ANES), and state-level hate crimes against Asians. Each panel presents state-level bias during a certain year. The boundaries in black represent the different Census divisions in the United States. Notice how there is a variation across states within a region.}}
\end{figure}
\end{center}

\newpage
\pagebreak

\begin{center}
\begin{figure}[H]
\caption{Maps of State-level Bias 2004--2021 Measure with Census Division Regional Boundaries}
\includegraphics[width=\textwidth]{Average_Skinmap.png} 
\label{fig:iat-map-all}
\caption*{\footnotesize{This figure shows the state-level bias index in the sample from 2004 to 2021. The bias units are in standard deviations and ranges from low to high bias. Bias index is constructed following \textcite{lubotskyInterpretationRegressionsMultiple2006}. The data is from the 2004--2021 Harvard's Project Implicit Association Test scores, American National Election Studies (ANES), and state-level hate crimes against Asians. The boundaries in black represent the different Census divisions in the United States. Notice how there is a variation across states within a region.}}
\end{figure}
\end{center}

\newpage
\pagebreak

% =========================
% Asian identity histograms
% for all generations
% =========================

\begin{center}
\begin{figure}[H]
\caption{Asian Racial Identity}
\includegraphics[width=\textwidth]{histogram_asian_american_race_aggregated.png} 
\label{fig:histogram-all}
\caption*{\footnotesize{This figure shows the distribution of Asian racial identity among respondents. 
The data is aggregated from the 2004--2021 Current Population Survey (CPS). 
The sample includes first-, second-, and third-generation objectively Asian Americans.
Following \textcite{antmanEthnicAttritionObserved2016,antmanEthnicAttritionAssimilation2020}, 
I utilize birthplace information for individuals, parents, and grandparents to create objective Asian ancestry indicators.
}}
\end{figure}
\end{center}

\newpage
\pagebreak

\begin{center}
\begin{figure}[H]
\caption{Asian Racial Identity: First Generation}
\includegraphics[width=\textwidth]{histogram_asian_american_race_firstgen.png} 
\label{fig:histogram-firstgen}
\caption*{\footnotesize{This figure shows the distribution of Asian racial identity among respondents. 
The data is aggregated from the 2004--2021 Current Population Survey (CPS). 
The sample includes first-generation objectively Asian Americans.
Following \textcite{antmanEthnicAttritionObserved2016,antmanEthnicAttritionAssimilation2020}, 
I utilize birthplace information for individuals, parents, and grandparents to create objective Asian ancestry indicators.
A first-generation Asian American is defined as an individual born in an Asian country and is not a US citizen born to US citizen parents abroad.
}}
\end{figure}
\end{center}


\pagebreak
\newpage

\begin{landscape}
\begin{figure}[!htb]
\centering
\caption{Asian Racial Identity: Second Generation}
\label{fig:histogram-secondgen}

% Top row
\begin{subfigure}{.48\textwidth}
\caption{All Second Generation}
\centering
\includegraphics[width=1\linewidth]{histogram_asian_american_race_secondgen.png}
\end{subfigure}
\hfill
\begin{subfigure}{.48\textwidth}
\caption{Second Generation: Asian Fathers-Asian Mothers}
\centering
\includegraphics[width=1\linewidth]{histogram_asian_american_race_secondgen_AA.png}
\end{subfigure}

\vspace{1cm}

% Bottom row
\begin{subfigure}{.48\textwidth}
\caption{Second Generation: Asian Fathers-White Mothers}
\centering
\includegraphics[width=1\linewidth]{histogram_asian_american_race_secondgen_AW.png}
\end{subfigure}
\hfill
\begin{subfigure}{.48\textwidth}
\caption{Second Generation: White Fathers-Asian Mothers}
\centering
\includegraphics[width=1\linewidth]{histogram_asian_american_race_secondgen_WA.png}
\end{subfigure}

\caption*{\footnotesize{This figure shows the distribution of Asian racial identity among respondents. 
The data is aggregated from the 2004--2021 Current Population Survey (CPS). 
The sample includes second-generation objectively Asian Americans.
Following \textcite{antmanEthnicAttritionObserved2016,antmanEthnicAttritionAssimilation2020}, 
I utilize birthplace information for individuals, parents, and grandparents to create objective Asian ancestry indicators.
A second-generation Asian American is defined as a native-born individual with at least one parent born in an Asian country. 
The first panel is for second-generation Asian Americans. The second panel is for second-generation Asian Americans with both parents born in an Asian country. 
The third panel is for second-generation Asian Americans with an Asian father and a White mother. 
The fourth panel is for second-generation Asian Americans with a White father and an Asian mother.
}}
\end{figure}
\end{landscape}

\pagebreak
\newpage

\pagebreak
\newpage

\begin{landscape}
\begin{figure}[!htb]
\centering
\caption{Asian Racial Identity: Third Generation}
\label{fig:histogram-thirdgen}

% Top row - 3 graphs
\begin{subfigure}{.32\textwidth}
\caption{All Third Generation}
\centering
\includegraphics[width=1\linewidth]{histogram_asian_american_race_thirdgen.png}
\end{subfigure}
\hfill
\begin{subfigure}{.32\textwidth}
\caption{Third Generation: One Asian Grandparent}
\centering
\includegraphics[width=1\linewidth]{histogram_asian_american_race_thirdgen_oneasiangran.png}
\end{subfigure}
\hfill
\begin{subfigure}{.32\textwidth}
\caption{Third Generation: Two Asian Grandparents}
\centering
\includegraphics[width=1\linewidth]{histogram_asian_american_race_thirdgen_twoasiangran.png}
\end{subfigure}

\vspace{1cm}

% Bottom row - 2 graphs centered with more space
\hspace*{\fill}
\begin{subfigure}{.4\textwidth}
\caption{Third Generation: Three Asian Grandparents}
\centering
\includegraphics[width=1\linewidth]{histogram_asian_american_race_thirdgen_threeasiangran.png}
\end{subfigure}
\hfill
\begin{subfigure}{.4\textwidth}
\caption{Third Generation: Four Asian Grandparents}
\centering
\includegraphics[width=1\linewidth]{histogram_asian_american_race_thirdgen_fourasiangran.png}
\end{subfigure}
\hspace*{\fill}

\caption*{\footnotesize{This figure shows the distribution of Asian racial identity among respondents. 
The data is aggregated from the 2004--2021 Current Population Survey (CPS). 
The sample includes second-generation objectively Asian Americans.
Following \textcite{antmanEthnicAttritionObserved2016,antmanEthnicAttritionAssimilation2020}, 
I utilize birthplace information for individuals, parents, and grandparents to create objective Asian ancestry indicators.
A third-generation Asian American is defined as a native-born individual with native-born parents and at least one grandparent born in an Asian country. 
The first panel is for third-generation Asian Americans. The second panel is for third-generation Asian Americans with one grandparent born in an Asian country. 
The third panel is for third-generation Asian Americans with two grandparents born in an Asian country. 
The fourth panel is for third-generation Asian Americans with three grandparents born in an Asian country.
The fifth panel is for third-generation Asian Americans with four grandparents born in an Asian country.
}}
\end{figure}
\end{landscape}

\pagebreak
\newpage

% =========================
% Multinomial Logit Marginal
% Effects Figures
% =========================

% Marginal Effects: Second Generation by Parental Type (AW and WA)
\begin{center}
\begin{figure}[!htb]
\centering
\caption{Marginal Effects of Key Covariates on Racial Identity Choice by Parental Ancestry (Second-Generation Asian Americans)}
\label{fig:marginal-effects-second-parental}

\begin{subfigure}{.48\textwidth}
\caption{Asian Father-White Mother (AW)}\label{subfig:meaw}
\centering
\includegraphics[width=1\linewidth]{boot_me_second_AW.png}
\end{subfigure}
\hfill
\begin{subfigure}{.48\textwidth}
\caption{White Father-Asian Mother (WA)}\label{subfig:mewa}
\centering
\includegraphics[width=1\linewidth]{boot_me_second_WA.png}
\end{subfigure}

\vspace{0.5cm}

\begin{subfigure}{.48\textwidth}
\caption{Pooled: All Mixed-Race}
\centering
\includegraphics[width=1\linewidth]{boot_me_second_AW_WA.png}
\end{subfigure}

\caption*{\footnotesize{This figure displays the marginal effects of key covariates on the probability of each racial identification outcome, estimated from equation (\ref{eq:multinomial_logit}) using multinomial logit models for second-generation Asian Americans by parental ancestry configuration. Marginal effects represent the change in predicted probability for a one-unit increase in the covariate (or discrete change from 0 to 1 for binary variables), holding other variables at their sample means. I include region $\times$ year fixed effects with controls for quartic age and local Asian population share. Each point corresponds to the median marginal effect across 1{,}000 bootstrap resamples, with 95\% percentile confidence intervals. Second-generation Asian Americans are native-born individuals with at least one Asian-born parent.}}
\end{figure}
\end{center}

\pagebreak
\newpage

% Marginal Effects: Second Generation by Parental Type (AW and WA)
\begin{center}
\begin{figure}[!htb]
\centering
\caption{Marginal Effects of Key Covariates on Racial Identity Choice by Parental Ancestry (Adult Second-Generation Asian Americans)}
\label{fig:marginal-effects-second-adults}

\begin{subfigure}{.48\textwidth}
\caption{Asian Father-White Mother (AW)}
\centering
\includegraphics[width=1\linewidth]{logit_boot_me_second_AW_adults_multinom.png}
\end{subfigure}
\hfill
\begin{subfigure}{.48\textwidth}
\caption{White Father-Asian Mother (WA)}
\centering
\includegraphics[width=1\linewidth]{logit_boot_me_second_WA_adults_multinom.png}
\end{subfigure}

\vspace{0.5cm}

\begin{subfigure}{.48\textwidth}
\caption{Pooled: All Mixed-Race Adults}
\centering
\includegraphics[width=1\linewidth]{logit_boot_me_second_AW_WA_adults_multinom.png}
\end{subfigure}

\caption*{\footnotesize{This figure displays the marginal effects of key covariates on the probability of each racial identification outcome, estimated from equation (\ref{eq:multinomial_logit}) using multinomial logit models for second-generation adults with Asian father/White mother (AW) and White father/Asian mother (WA). Marginal effects represent the change in predicted probability for a one-unit increase in the covariate (or discrete change from 0 to 1 for binary variables), holding other variables at their sample means. I include region $\times$ year fixed effects with controls for quartic age and local Asian population share. Each point corresponds to the median marginal effect across 1{,}000 bootstrap resamples, with 95\% percentile confidence intervals. Second-generation Asian Americans are native-born individuals with at least one Asian-born parent.}}
\end{figure}
\end{center}

\pagebreak
\newpage

\begin{center}
\begin{figure}[!htb]
\centering
\caption{Marginal Effects of Key Covariates on Racial Identity Choice by Number of Asian Grandparents (Third-Generation Asian Americans)}
\label{fig:marginal-effects-third-grandparental}
\begin{subfigure}{.45\textwidth}
\caption{One Asian Grandparent}\label{subfig:oneasiangrand}
\centering
\includegraphics[width=1\linewidth]{boot_me_third_oneAsian.png}
\end{subfigure}
\hfill
\begin{subfigure}{.45\textwidth}
\caption{Two Asian Grandparents}\label{subfig:twoasiangrand}
\centering
\includegraphics[width=1\linewidth]{boot_me_third_twoAsian.png}
\end{subfigure}

\vspace{0.5cm}

\begin{subfigure}{.45\textwidth}
\caption{Three Asian Grandparents}
\centering
\includegraphics[width=1\linewidth]{boot_me_third_threeAsian.png}
\end{subfigure}
\hfill
\begin{subfigure}{.45\textwidth}
\caption{Pooled: One and Two Asian Grandparents}
\centering
\includegraphics[width=1\linewidth]{boot_me_third_oneTwoAsian.png}
\end{subfigure}

\vspace{0.5cm}

\begin{subfigure}{.45\textwidth}
\caption{Pooled: All Third Generation}
\centering
\includegraphics[width=1\linewidth]{boot_me_third_oneTwoThreeAsian.png}
\end{subfigure}

\caption*{\footnotesize{This figure displays the marginal effects of key covariates on the probability of each racial identification outcome, estimated from equation (\ref{eq:multinomial_logit}) using multinomial logit models for third-generation Asian Americans by number of Asian grandparents. Marginal effects represent the change in predicted probability for a one-unit increase in the covariate (or discrete change from 0 to 1 for binary variables), holding other variables at their sample means. I include region $\times$ year fixed effects with controls for quartic age and local Asian population share. Each estimate is the median across 1{,}000 bootstrap resamples, and 95\% percentile confidence intervals accompany every point.}}
\end{figure}
\end{center}

\pagebreak
\newpage

% =========================
% 42: Interaction Coefficient Plots (AW/WA)
% =========================
\begin{center}
\begin{figure}[!htb]
\centering
\caption{Interaction Effects of Anti-Asian Bias with Gender and Parental Education: By Mixed Parental Types Among Adults}
\label{fig:interaction-coefs-aw-wa}

\begin{subfigure}{.48\textwidth}
\caption{Asian Fathers-White Mothers}
\centering
\includegraphics[width=1\linewidth]{interaction_coefficients_AW.png}
\end{subfigure}
\hfill
\begin{subfigure}{.48\textwidth}
\caption{White Fathers-Asian Mothers}
\centering
\includegraphics[width=1\linewidth]{interaction_coefficients_WA.png}
\end{subfigure}

\vspace{0.5cm}

\begin{subfigure}{.48\textwidth}
\caption{Pooled: All Mixed-Race Adults}
\centering
\includegraphics[width=1\linewidth]{interaction_coefficients_AWWA_adults.png}
\end{subfigure}

\caption*{\footnotesize{I show interaction coefficients from estimating $A_{ist}^g = \beta_1^g AntiAsianBias_{st} + \mathbf{X}_{ist}^g\pi + \gamma_{rt} + \varepsilon_{ist}$ with interaction terms for gender and parental education. I include region $\times$ year fixed effects with controls for sex, quartic age, years of education, and the inverse hyperbolic sine of income. The dependent variable is self-reported Asian identity and the independent variables include demeaned state-level bias interacted with gender and parental education indicators. Each panel results from separate regressions on different samples divided by mixed parental types. Standard errors are clustered on the state level. The samples include second-generation Asian individuals ages 18 and above with mixed-race parentage. Native-born second-generation Asian immigrant individuals with one parent born in an Asian country and one parent born outside of Asia.}}
\end{figure}
\end{center}

\pagebreak
\newpage

% =========================
% 44: Interaction Coefficient Plots (AW/WA/AWWA - All Ages)
% =========================
\begin{center}
\begin{figure}[!htb]
\centering
\caption{Interaction Effects of Anti-Asian Bias with Gender and Parental Education: By Mixed Parental Types}
\label{fig:interaction-coefs-aw-wa-awwa}

\begin{subfigure}{.45\textwidth}
\caption{Asian Fathers-White Mothers}
\centering
\includegraphics[width=1\linewidth]{interaction_coefficients_AW.png}
\end{subfigure}
\hfill
\begin{subfigure}{.45\textwidth}
\caption{White Fathers-Asian Mothers}
\centering
\includegraphics[width=1\linewidth]{interaction_coefficients_WA.png}
\end{subfigure}

\vspace{0.5cm}

\begin{subfigure}{.45\textwidth}
\caption{Pooled: All Mixed-Race}
\centering
\includegraphics[width=1\linewidth]{interaction_coefficients_AWWA.png}
\end{subfigure}

\caption*{\footnotesize{I show interaction coefficients from estimating $A_{ist}^g = \beta_1^g AntiAsianBias_{st} + \mathbf{X}_{ist}^g\pi + \gamma_{rt} + \varepsilon_{ist}$ with interaction terms for gender and parental education. I include region $\times$ year fixed effects with controls for sex, quartic age, years of education, and the inverse hyperbolic sine of income. The dependent variable is self-reported Asian identity and the independent variables include demeaned state-level bias interacted with gender and parental education indicators. Each panel results from separate regressions on different samples divided by mixed parental types. Standard errors are clustered on the state level. The samples include second-generation Asian individuals with mixed-race parentage. Native-born second-generation Asian immigrant individuals with one parent born in an Asian country and one parent born outside of Asia.}}
\end{figure}
\end{center}

\pagebreak
\newpage

% =========================
% 43: Interaction Coefficient Plots (Third Generation by Grandparent Type)
% =========================
\begin{center}
\begin{figure}[!htb]
\centering
\caption{Interaction Effects of Anti-Asian Bias with Gender and Parental Education: By Asian Grandparent Count Among Third Generation}
\label{fig:interaction-coefs-thirdgen-grandparent}

\begin{subfigure}{.45\textwidth}
\caption{One Asian Grandparent}
\centering
\includegraphics[width=1\linewidth]{interaction_coefficients_OneAsian.png}
\end{subfigure}
\hfill
\begin{subfigure}{.45\textwidth}
\caption{Two Asian Grandparents}
\centering
\includegraphics[width=1\linewidth]{interaction_coefficients_TwoAsian.png}
\end{subfigure}

\vspace{0.5cm}

\begin{subfigure}{.45\textwidth}
\caption{Three Asian Grandparents}
\centering
\includegraphics[width=1\linewidth]{interaction_coefficients_ThreeAsian.png}
\end{subfigure}
\hfill
\begin{subfigure}{.45\textwidth}
\caption{Pooled: All Third Generation}
\centering
\includegraphics[width=1\linewidth]{interaction_coefficients_pooled_thirdgen.png}
\end{subfigure}

\caption*{\footnotesize{I show interaction coefficients from estimating $A_{ist}^g = \beta_1^g AntiAsianBias_{st} + \mathbf{X}_{ist}^g\pi + \gamma_{rt} + \varepsilon_{ist}$ with interaction terms for gender and parental education on third-generation samples. I include region $\times$ year fixed effects with controls for sex, quartic age, years of education, and the inverse hyperbolic sine of income. The dependent variable is self-reported Asian identity and the independent variables include demeaned state-level bias interacted with gender and parental education indicators. Each panel results from separate regressions on different samples divided by the number of Asian grandparents. Standard errors are clustered on the state level. The samples include third-generation Asian individuals with varying degrees of Asian ancestry through grandparents. Native-born third-generation individuals with at least one grandparent born in an Asian country.}}
\end{figure}
\end{center}